\documentclass[10pt,a4paper]{amsart}
\usepackage[margin=19mm]{geometry}
\usepackage{amsmath,amssymb,mathtools}
\usepackage[scr=rsfs,cal=euler]{mathalfa}
\usepackage{xfrac}
\usepackage[unicode,hidelinks,bookmarks=false]{hyperref}
\newcommand{\ie}{i.e.\ }
\newcommand{\cf}{cf.\ }
\newcommand{\resp}{respectively\ }
\newcommand{\Subsection}{\subsection}
\providecommand{\nobreakdash}{\nobreak\hbox{-}}
\setlength{\emergencystretch}{2em}
\multlinegap0pt
\usepackage{amssymb}
\usepackage[shortlabels]{enumitem}
\newtheorem{thm}{Theorem}
\newtheorem{lem}{Lemma}
\title{A further $q$-generalization of the (C.2) and (G.2) supercongruences of Van Hamme}
\author{Song-Xiao Li, Su-Dan Wang}
\date{Source: arXiv:2603.26223v1 (CC BY 4.0)}
\begin{document}
\maketitle

\noindent\emph{Source and licence.} A further $q$-generalization of the (C.2) and (G.2) supercongruences of Van Hamme. Version arXiv:2603.26223v1 (CC BY 4.0), distributed by arXiv under the Creative Commons Attribution 4.0 International licence, http://creativecommons.org/licenses/by/4.0/, as recorded in the arXiv licence metadata for this exact version and shown on the abstract page https://arxiv.org/abs/2603.26223v1. Full licence notice: \texttt{licenses/arxiv-2603.26223-CC-BY-4.0.txt}.\par\smallskip

\noindent\textit{Editorial scope.} Counted unit: the article's thm thm:main-2 (source0.tex lines 194--204). The article's complete original proof of it is delivered with it (source0.tex lines 451--508, one \texttt{proof} environment of 58 lines, which the article places away from the statement). No mathematics is written, repaired or re-derived: the statement and the proof are the article's own lines, in the article's own notation, and no retained line is substituted or re-flowed.  Retained uncounted as the source-local context the proof needs: lines 357--359; lines 411--414; lines 428--436; lines 439--443; lines 445--449.  Nothing was omitted from any retained span, so the package carries no omission record.  No retained line contains a citation, so the wrapper carries no bibliography and no bibliography entry.  No retained line contains a figure environment, an image-inclusion command or drawing code, so no image is delivered and the package declares no asset dependency.  Editorial formatting only, in addition: an AMS \texttt{amsart} preamble that restates the article's own declarations and macros used by the retained lines, namely the environments \texttt{thm}, \texttt{lem} on separate counters, together with the standard packages those lines need.  The retained environments print in this document as Theorem 1, Lemma 1 in order of appearance, whereas the article prints the same items with the numbering of its own full document.  The complete native source of the article as distributed in the arXiv e-print for this exact version is bundled unchanged as \texttt{sources/197267/source0.tex}.
\begin{align}\label{eq:qz}
    [dk-d+r]F(m,k-1)-[dk-d+2r]F(m,k)=G(m+1,k)-G(m,k).
\end{align}

\begin{align}\label{eq:iden}
    \frac{(q^d;q^d)_{k}(q^{2r};q^d)_{k-1}}{[dk][dk-d+r]^2 (q^r;q^d)_k(q^r;q^d)_{k-1}}=
    \frac{(q^d;q^d)_{k-1}(q^{2r};q^d)_{k-1}}{[dk-d+r]^3(q^r;q^d)^2_{k-1}},
\end{align}

\begin{lem}\label{lem:F-dn}
    Let $n$ and $d\geqslant2$ be positive integers satisfy $\gcd(n,d)=1$. Let $r$ be an integer with $d-n\leqslant r\leqslant (d-1)n$, and $n\equiv -r\pmod{d}$.
    Then modulo $[n]\Phi_n(q)^4$,
    \begin{align*}\label{eq:F-dn}
        &q^{((d-1)n-r)(d-(d-1)n-3r)/2d}[2(d-1)n-r]\frac{(q^r;q^d)_{2((d-1)n-r)/d}}{(q^d;q^d)^2_{((d-1)n-r)/d}}\nonumber\\
        &\hspace{2em}\equiv(-1)^{((d-1)n-r)/d} q^{r(n+r-dn)/d}\nonumber\\
        &\hspace{3em}\times\left\{[(d-1)n]-([(d-1)n]^3+[(d-1)n]^4(1-q))\sum_{k=1}^{((d-1)n-r)/d}\frac{q^{kd}}{[kd]^2}\right\}.
    \end{align*}
\end{lem}

\begin{lem}\label{lem:G-div-dn}
    Let $n$ and $d\geqslant2$ be positive integers satisfy $\gcd(n,d)=1$. Let $r$ be an integer with $d-n\leqslant r\leqslant (d-1)n$, and $n\equiv -r\pmod{d}$.
    For $k\in \{1,2,\dots, \frac{(d-1)n-r}{d}\}$, write \[\frac{(1-q^{(d-1)n})(q^d;q^d)_{k-1}(q^{2r};q^d)_{k-1}}{(q^r;q^d)_{k}(q^r;q^d)_{k-1}}
    =\frac{A_k(q)}{B_k(q)},\] where $A_k(q)$ and $B_k(q)$ are relatively prime polynomials in $q$. Then $B_k(q)$ is relatively prime to $1-q^n$.
\end{lem}

\begin{lem}\label{lem:q-integers-dn}
    Let $n$ and $d\geqslant2$ be positive integers satisfy $\gcd(n,d)=1$. Let $r$ be an integer with $d-n\leqslant r\leqslant (d-1)n$, and $n\equiv -r\pmod{d}$.
    For $k\in \{1,2,\dots, \frac{(d-1)n-r}{d}\}$, \[\frac{1-q^{(d-1)n}}{1-q^{dk}}=\frac{C_k(q)}{D_k(q)},\] where $C_k(q)$ and $D_k(q)$ are relatively prime polynomials in $q$. 
    Then $D_k(q)$ is relatively prime to $1-q^n$.
\end{lem}

\begin{thm}\label{thm:main-2}
    Let $n$ and $d\geqslant2$ be positive integers satisfying $\gcd(n,d)=1$. Let $r$ be an integer with $d-n\leqslant r \leqslant dn-n$ and 
    $n\equiv -r \pmod{d} $. Then, modulo $[n]\Phi_{n}(q)^4$,
    \begin{align*}
        &\sum_{k=0}^{M}[2dk+r]\frac{(q^r;q^d)_k^4}{(q^d;q^d)_k^4}q^{(d-2r)k}\nonumber\\
        &\hspace{2em}\equiv q^{r(n+r-dn)/d}\frac{(q^{2r};q^d)_{((d-1)n-r)/d}}{(q^{d};q^d)_{((d-1)n-r)/d}}\nonumber\\
        &\hspace{3em}\times\Bigg([(d-1)n]-([(d-1)n]^3+[(d-1)n]^4(1-q))\sum_{k=1}^{((d-1)n-r)/d}\frac{q^{kd}}{[kd]^2}\Bigg)\nonumber\\
        &\hspace{3em}+[(d-1)n]^4\sum_{k=1}^{((d-1)n-r)/d}q^{(kd-d-2dn+2n+2r)(d-dn+n)/d}\frac{(q^d;q^d)_{k-1}(q^{2r};q^d)_{k-1}}{[dk-d+r]^3(q^r;q^d)_{k-1}^2},
    \end{align*}
    where $M=((d-1)n-r)/d$ or $n-1$.
\end{thm}

\begin{proof}[Proof of Theorem~\ref{thm:main-2}]
Let \(F(m,k)\) and \(G(m,k)\) be as defined in Section~3.
Summing \eqref{eq:qz} with respect to \(m\) from \(0\) to \(((d-1)n-r)/d\), we obtain
\begin{align*}
    G&\left(\frac{(d-1)n+d-r}{d},k\right)\\
    &=\sum_{m=0}^{((d-1)n-r)/d}[dk-d+r]F(m,k-1)-\sum_{m=0}^{((d-1)n-r)/d}[dk-d+2r]F(m,k),
\end{align*}
where we have used the fact that $G(0,k)=0$.
Then, summing over $k$ from $1$ to $((d-1)n-r)/d$ and carrying out an iterative computation,
 observing that $F(m,((d-1)n-r)/d)=0$ for $m<((d-1)n-r)/d$, we obtain
\begin{align}\label{eq:FG-dn}
    \sum_{m=0}^{((d-1)n-r)/d}F(m,0)
    &=F\left(\frac{(d-1)n-r}{d},\frac{(d-1)n-r}{d}\right)\prod_{j=1}^{((d-1)n-r)/d}\frac{[dj-d+2r]}{[dj-d+r]}\nonumber\\
    &\hspace{1em}+\sum_{k=1}^{((d-1)n-r)/d}G\left(\frac{(d-1)n+d-r}{d},k\right)\prod_{j=1}^{k-1}\frac{[dj-d+2r]}{[dj-d+r]}\frac{1}{[dk-d+r]}.
\end{align}
By Lemma~\ref{lem:F-dn}, we conclude that, modulo $[n]\Phi_n(q)^4$,
\begin{align}\label{eq:F-dn-final}
    &F\left(\frac{(d-1)n-r}{d},\frac{(d-1)n-r}{d}\right)\prod_{j=1}^{((d-1)n-r)/d}\frac{[dj-d+2r]}{[dj-d+r]}\nonumber\\
    &\hspace{1em}=(-1)^{((d-1)n-r)/d}q^{(dn-n-r)(d-dn+n-3r)/2d}[2(d-1)n-r]\nonumber\\
    &\hspace{2em}\times\frac{(q^r;q^d)_{2((d-1)n-r)/d}}{(q^d;q^d)^2_{((d-1)n-r)/d}}\cdot
    \frac{(q^{2r};q^d)_{((d-1)n-r)/d}}{(q^d;q^d)_{((d-1)n-r)/d}}\nonumber\\
    &\hspace{1em}\equiv q^{r(r-(d-1)n)/d}\frac{(q^{2r};q^d)_{((d-1)n-r)/d}}{(q^d;q^d)_{((d-1)n-r)/d}}\nonumber\\
    &\hspace{2em}\times\left\{[(d-1)n]-([(d-1)n]^3+[(d-1)n]^4(1-q))
    \sum_{k=1}^{((d-1)n-r)/d}\frac{q^{kd}}{[kd]^2}\right\}.
\end{align}
For $k=1,2,\dots,((d-1)n-r)/d$, we see that 
\begin{align}\label{eq:G-dn}
    G\left(\frac{(d-1)n+d-r}{d},k\right)&=(-1)^{k-1}q^{d\binom{k}{2}-((d-1)n-r)(k-(d-2r)/d)}\nonumber\\
    &\hspace{1em}\times\frac{(q^r;q^d)^3_{((d-1)n-r)/d+1}(q^r;q^d)_{((d-1)n-r)/d+k}}
    {(1-q)^2(q^d;q^d)^3_{((d-1)n-r)/d}(q^d;q^d)_{((d-1)n-r)/d-k+1}(q^r;q^d)_k^2}\nonumber\\
    &\hspace{0.5em}=(-1)^{k-1}q^{d\binom{k}{2}-((d-1)n-r)(k-(d-2r)/d)}\nonumber\\
    &\hspace{1em}\times\frac{(1-q^{(d-1)n})^4(q^r;q^d)^4_{((d-1)n-r)/d}(q^{(d-1)n+d};q^d)_{k-1}}
    {(1-q)^2(q^d;q^d)^3_{((d-1)n-r)/d}(q^d;q^d)_{((d-1)n-r)/d-k+1}(q^r;q^d)^2_k}.
\end{align}
Making use of the fact that $q^n\equiv 1 \pmod{\Phi_n(q)}$, it is readily seen that
\begin{align*}
    (q^d;q^d)_{((d-1)n-r)/d-k+1}\equiv (-1)^{k-1}q^{(k-1)(dk+2r-2d)/2}\frac{(q^d;q^d)_{((d-1)n-r)/d}}{(q^r;q^d)_{k-1}}\pmod{\Phi_n(q)}
\end{align*}
and 
\begin{align*}
    \frac{(q^r;q^d)_{((d-1)n-r)/d}}{(q^d;q^d)_{((d-1)n-r)/d}}\equiv (-1)^{((d-1)n-r)/d}q^{(dn-n-r)(dn-n+r-d)/2d}\pmod{\Phi_n(q)}.
\end{align*}
In view of the above two $q$-congruences, it follows from \eqref{eq:G-dn} that, for $1\leqslant k\leqslant((d-1)n-r)/d$,
\begin{align*}
    G&\left(\frac{(d-1)n+d-r}{d},k\right)\\
    &\equiv q^{(kd-d-2dn+2n+2r)(d-dn+n)/d}\frac{[(d-1)n]^4(q^d;q^d)_{k}}{[dk][dk-d+r](q^r;q^d)_{k}}\pmod{\Phi_n(q)^5}.
\end{align*}
Furthermore, we obtain
\begin{align}\label{eq:G-dn-final}
    &G\left(\frac{(d-1)n+d-r}{d},k\right)\prod_{j=1}^{k-1}\frac{[dj-d+2r]}{[dj-d+r]}\cdot\frac{1}{[dk-d+r]}\nonumber\\
    &\hspace{1em}\equiv q^{(kd-d-2dn+2n+2r)(d-dn+n)/d}\frac{[(d-1)n]^4(q^d;q^d)_{k}(q^{2r};q^d)_{k-1}}{[dk][dk-d+r]^2(q^r;q^d)_{k}(q^r;q^d)_{k-1}}\pmod{\Phi_n(q)^5}.
\end{align}
By virtue of Lemma~\ref{lem:G-div-dn} and Lemma~\ref{lem:q-integers-dn}, it is easy to see that the right-hand side 
of \eqref{eq:G-dn-final} is congruent to 0 modulo $[n]$. Similar as before, we can obtain that $G(\tfrac{(d-1)n+d-r}{d},k)\equiv 0\pmod{[n]}$ from 
\eqref{eq:G-dn}, and so is the left-hand side of \eqref{eq:G-dn-final}. Namely, the $q$-congruence \eqref{eq:G-dn-final} holds modulo $[n]$.
It follows from the identity $\operatorname{lcm}([n],\Phi_n(q)^5)=[n]\Phi_n(q)^4$ that the $q$-congruence \eqref{eq:G-dn-final} holds modulo $[n]\Phi_n(q)^4$.\par
Combining \eqref{eq:FG-dn}, \eqref{eq:F-dn-final} and the case modulo $[n]\Phi_n(q)^4$, we immediately complete the proof by using the identity \eqref{eq:iden}.
\end{proof}

\end{document}
